\chapter{Filosofía de la protección radiológica}

Tanto en lo que se refiere a trabajadores ocupacionalmente expuestos como al público, se deberá procurar que las dosis se mantengan en un nivel mínimo, tomando en cuenta consideraciones de riesgo-beneficio.

Esto es contemplado por el criterio ``ALARA'' (\emph{As Low As is Reasonably Achievable}), que significa «tan bajo como razonablemente pueda alcanzarse».

Todas las actividades que involucren exposición a la radiación deberán estar gobernadas por este criterio, y la evaluación ALARA de una cierta actividad podrá servir como base en la toma de decisiones para su modificación o corrección.

Debido a que en las actividades normales de la radiografía industrial es frecuente trabajar en áreas imposibles de evacuar, es necesario tomar una serie de medidas tendientes a proteger al público o a trabajadores no ocupacionalmente expuestos.

Dentro de estas medidas se incluyen la delimitación de áreas y su señalización adecuada. Principalmente se busca definir dos tipos de zonas de acuerdo con lo siguiente.

\section{Zona restringida (zona de radiación)}

Es aquella a la que, por motivos de seguridad radiológica, se prohíbe el acceso o permanencia a todo el personal, excepto al técnico encargado de la fuente o a sus auxiliares y siempre que dicha fuente no se encuentre fuera de su contenedor.

Toda zona con niveles de exposición iguales o superiores a $5\ \mathrm{mR/h}$ será considerada como zona restringida. La zona restringida se mantendrá acordonada y debidamente señalizada con letreros que indiquen la presencia de radiación, y con distancias entre dos letreros consecutivos no mayores a 10 metros.

\section{Zona controlada}

En las áreas donde el nivel de radiactividad es controlado, se vigila el flujo de personal. Se considera como zona controlada aquella comprendida fuera de la restringida pero con un nivel de exposición mayor a $2\ \mathrm{mR/h}$. La zona controlada no requiere acordonamiento, pero se mantendrá vigilada por el personal responsable, a fin de evitar que alguna persona pudiera permanecer innecesariamente en ese lugar.

\section{Límites administrativos de dosis}

Los límites administrativos de dosis de radiación serán los indicados como límites operativos de IINDESA en nuestro Manual de Seguridad Radiológica:
\begin{itemize}
  \item El P.O.E. no deberá recibir una dosis mayor a $4\ \mathrm{mSv}$ ($400\ \mathrm{mrem}$) por mes.
  \item El P.O.E. no deberá rebasar $5\ \mathrm{mSv}$ ($5000\ \mathrm{mrem}$) en un año.
\end{itemize}

Se deberá cumplir con los niveles de referencia y de registro indicados en nuestro Manual de Seguridad Radiológica, parte 3.2.

Las restricciones sobre la dosis efectiva son suficientes para asegurar que se eviten los efectos determinísticos en todos los tejidos y órganos del cuerpo, excepto el cristalino del ojo, el cual hace una contribución despreciable a la dosis efectiva, y la piel, que puede estar sujeta a exposiciones localizadas. Se necesitan límites especiales para estos tejidos. Los límites anuales son $150\ \mathrm{mSv}$ para el cristalino y $500\ \mathrm{mSv}$ para la piel, promediados sobre cualquier centímetro cuadrado, independientemente del área expuesta.
